% Section 4.4, from lectures/L04/appendix/c_exercises.md.

\section{Exercises}
\label{c4:sec:exercises}\label{c4:app:c}
These exercises reinforce the concepts from \secref{c4:app:a} and \secref{c4:app:b}.

\begin{aside}
\textbf{How to check your work.} Write each exercise set as its own application under
\file{lectures/L04/exercises}, with the file names given there, and run \code{make test} from the
repository root. A set’s tests switch on as soon as its files exist, and a set you have not started
is reported as \code{SKIP}. Each application grows through its set, and the tests check where it
ends up. Programs run on \code{native_sim} in simulated time, so their output, timestamps included,
is compared with the expected output line by line. The threads’ priorities, start delays and sleeps
given below are part of the exercise: change one, and the output changes.

The solutions are published in the companion repository, under
\file{lectures/L04/appendix/solutions}, a few days after the lecture. Try each exercise yourself
before you read them.
\end{aside}

Remember \code{source env.sh} in every new terminal. Build and run an exercise from the repository
root:

\begin{shell}
west build -p -b native_sim/native/64 lectures/L04/exercises/exercise1
./build/zephyr/zephyr.exe -stop_at=3 -no-rt
\end{shell}

% ---------------------------------------------------------------------
\exerciseset{Readings in a queue}
\label{c4:set:1}
Create an application named \code{exercise1} to work in for this exercise set, with a
\code{CMakeLists.txt}, a \code{prj.conf} and a \code{src/main.c}.

\exercise{A sampler and a slower storage thread}{Code}
\label{c4:ex:1.1}
In this exercise you will connect a fast producer and a slow consumer with a message queue.

\task{Tasks}
\begin{itemize}
  \item In \code{prj.conf}, set one tick per millisecond, as this chapter’s examples do:
    \code{CONFIG_SYS_CLOCK_TICKS_PER_SEC=1000}.
  \item Define a message type:
\begin{ccode}
struct reading {
  uint32_t seq;
  int64_t timestamp_ms;
};
\end{ccode}
  \item Define a message queue named \code{readings} with room for \textbf{3} readings.
  \item Write a thread \code{sampler}, priority \textbf{5}, which takes 8 readings, numbered 1 to 8. For
    each:
\begin{itemize}
  \item Fill in a \code{struct reading} with the number and \code{k_uptime_get()}.
  \item Put it into the queue \textbf{without waiting}. If that fails, print \code{[<t> ms] queue full,
    dropped \#<n>}, with the reading’s number. Otherwise print \code{[<t> ms] sampled \#<n>, queue
    <used>/3}.
  \item Sleep 100 ms.
\end{itemize}
After the eighth reading, print \code{[<t> ms] sampler done}.
  \item Write a thread \code{storage}, priority \textbf{6}, which forever takes a reading from the queue,
    waiting as long as it takes, prints \code{[<t> ms]   stored \#<n>, <age> ms old} (three spaces after
    the timestamp), where the age is the uptime minus the reading’s timestamp, and sleeps 350 ms.
  \item Start both threads with \code{K_THREAD_DEFINE()}, a stack of 1024 bytes, and no start delay.
\end{itemize}

Print every timestamp with \code{\%4lld}. With \code{-stop_at=3}, the output must be:

\begin{console}
[   0 ms] sampled #1, queue 1/3
[   0 ms]   stored #1, 0 ms old
[ 101 ms] sampled #2, queue 1/3
[ 202 ms] sampled #3, queue 2/3
[ 303 ms] sampled #4, queue 3/3
[ 351 ms]   stored #2, 250 ms old
[ 404 ms] sampled #5, queue 3/3
[ 505 ms] queue full, dropped #6
[ 606 ms] queue full, dropped #7
[ 702 ms]   stored #3, 500 ms old
[ 707 ms] sampled #8, queue 3/3
[ 808 ms] sampler done
[1053 ms]   stored #4, 750 ms old
[1404 ms]   stored #5, 1000 ms old
[1755 ms]   stored #8, 1048 ms old
\end{console}

\noindent \textbf{Hint:} \code{K_MSGQ_DEFINE()}’s last argument is the message’s alignment:
\code{alignof(struct reading)}, from \code{<stdalign.h>}. \code{k_msgq_num_used_get()} returns how
many messages are queued.

\task{Reflection}
\begin{itemize}
  \item \code{sampled \#1} is printed before \code{stored \#1}. Why? Give \code{storage} priority 4 instead,
    and run it again: what changes, and why does the queue read \code{0/3}?
  \item Why does the sampler put without waiting, rather than with \code{K_FOREVER}? Which kind of code
    could not wait at all?
\end{itemize}

\exercise{Keep the newest}{Code}
\label{c4:ex:1.2}

\task{Tasks}
Change the sampler so that a full queue drops the \textbf{oldest} reading instead of the new one:

\begin{itemize}
  \item When the put fails, take the oldest reading out of the queue without waiting, print \code{[<t> ms]
    queue full, dropped \#<n>} with \textbf{its} number, and put the new reading again.
  \item Then print \code{sampled \#<n>} for the new reading, as for any other.
\end{itemize}

The output must now be:

\begin{console}
[   0 ms] sampled #1, queue 1/3
[   0 ms]   stored #1, 0 ms old
[ 101 ms] sampled #2, queue 1/3
[ 202 ms] sampled #3, queue 2/3
[ 303 ms] sampled #4, queue 3/3
[ 351 ms]   stored #2, 250 ms old
[ 404 ms] sampled #5, queue 3/3
[ 505 ms] queue full, dropped #3
[ 505 ms] sampled #6, queue 3/3
[ 606 ms] queue full, dropped #4
[ 606 ms] sampled #7, queue 3/3
[ 702 ms]   stored #5, 298 ms old
[ 707 ms] sampled #8, queue 3/3
[ 808 ms] sampler done
[1053 ms]   stored #6, 548 ms old
[1404 ms]   stored #7, 798 ms old
[1755 ms]   stored #8, 1048 ms old
\end{console}

\task{Reflection}
\begin{itemize}
  \item Compare the readings stored, and their ages, with Exercise~\ref{c4:ex:1.1}. Which policy does a
    temperature logger want? Which would a queue of commands want?
  \item Make the queue 8 readings long. Is anything dropped now? How old is the last reading when it is
    stored? Did the bigger queue help?
  \item How much RAM does the queue’s buffer use on \code{native_sim/native/64}, and when is it allocated?
  \item Dropping the oldest takes two calls, a get and a put. Why is that safe in this program, and when
    would it not be?
\end{itemize}

% ---------------------------------------------------------------------
\exerciseset{Priority inversion}
\label{c4:set:2}
Create an application named \code{exercise2} to work in for this exercise set.

\exercise{A log behind a semaphore}{Code}
\label{c4:ex:2.1}
In this exercise you will produce priority inversion, measure it, and fix it.

\task{Tasks}
\begin{itemize}
  \item In \code{prj.conf}, set \code{CONFIG_SYS_CLOCK_TICKS_PER_SEC=1000}.
  \item Create a \code{Kconfig} file with a menu named \code{Application} holding a \code{bool} option
    \code{APP_USE_MUTEX}, with the prompt \code{Protect the log with a k_mutex instead of a k_sem}, and
    end it so that Zephyr’s own options are still available. The option stays unused until
    Exercise~\ref{c4:ex:2.2}.
  \item In \code{src/main.c}, protect a shared log with a binary semaphore named \code{log_lock}, created
    available (initial count 1, limit 1).
  \item Write a function \code{trace(who, what)} that prints \code{[<t> ms] <who> (priority <p>) <what>},
    with the uptime as \code{\%4lld}, \code{who} left-aligned in 7 characters (\code{\%-7s}), and the
    calling thread’s current priority from \code{k_thread_priority_get()}.
  \item Write a function \code{compute(ms)} that uses CPU time by calling \code{k_busy_wait(USEC_PER_MSEC)}
    \code{ms} times.
  \item Write three threads, all with a stack of 1024 bytes:
\begin{itemize}
  \item \code{logger}, priority \textbf{9}, no start delay: traces \code{takes the lock}, takes it, computes
    for \textbf{60} ms, traces \code{gives the lock back}, and gives it back.
  \item \code{control}, priority \textbf{5}, start delay \textbf{10} ms: records the uptime, traces
    \code{needs the lock}, takes it, traces \code{has the lock}, gives it back, and prints \code{control
    waited <ms> ms}, the time since the uptime it recorded.
  \item \code{network}, priority \textbf{7}, start delay \textbf{20} ms: traces \code{starts sending},
    computes for \textbf{300} ms, and traces \code{is done}.
\end{itemize}
\end{itemize}

The output must be:

\begin{console}
[   0 ms] logger  (priority 9) takes the lock
[  11 ms] control (priority 5) needs the lock
[  21 ms] network (priority 7) starts sending
[ 321 ms] network (priority 7) is done
[ 360 ms] logger  (priority 9) gives the lock back
[ 360 ms] control (priority 5) has the lock
control waited 349 ms
\end{console}

\noindent \textbf{Hint:} why \code{compute()} busy-waits one millisecond at a time rather than
calling \code{k_busy_wait(60000)} once is in \secref{c4:app:a}, \emph{On \code{native_sim}: making a
thread use CPU time}.

\task{Reflection}
\begin{itemize}
  \item Before you run it, predict how long \code{control} waits. Then account for the 349 ms, millisecond
    by millisecond: who ran when?
  \item \code{network} never touches the lock. How did it delay \code{control} by 300 ms?
\end{itemize}

\exercise{The mutex}{Code}
\label{c4:ex:2.2}

\task{Tasks}
\begin{itemize}
  \item When \code{CONFIG_APP_USE_MUTEX} is enabled, protect the log with a \code{k_mutex} named
    \code{log_lock} instead of the semaphore, and change nothing else. Hide the difference behind two
    macros, such as \code{lock()} and \code{unlock()}, so that the threads’ code is the same in both
    builds.
  \item Create a file \code{mutex.conf} that enables the option.
\end{itemize}

Built with \code{-DEXTRA_CONF_FILE=mutex.conf}, the output must be:

\begin{console}
[   0 ms] logger  (priority 9) takes the lock
[  11 ms] control (priority 5) needs the lock
[  60 ms] logger  (priority 5) gives the lock back
[  60 ms] control (priority 5) has the lock
control waited 49 ms
[  60 ms] network (priority 7) starts sending
[ 360 ms] network (priority 7) is done
\end{console}

\noindent Built without it, the output must be unchanged from Exercise~\ref{c4:ex:2.1}.

\task{Reflection}
\begin{itemize}
  \item Why is \code{logger}’s priority 5 when it gives the lock back? What is it after the unlock?
  \item \code{network} was ready at 21 ms. Why does it only start sending at 60 ms?
  \item \code{control} still waited 49 ms. Could any lock bring that down to zero? What would?
  \item With the mutex, make \code{control} wait at most 20 ms, with \code{k_mutex_lock(&log_lock, K_MSEC(20))},
and trace \code{gave up} when that fails. What happens to \code{logger}’s priority when
\code{control} gives up, and what does \code{network} do then?
\end{itemize}

% ---------------------------------------------------------------------
\exerciseset{A battery on zbus}
\label{c4:set:3}
Create an application named \code{exercise3} to work in for this exercise set.

\exercise{The channel and a listener}{Code}
\label{c4:ex:3.1}
In this exercise you will publish battery readings on a zbus channel, and react to them.

\task{Tasks}
\begin{itemize}
  \item Enable zbus in \code{prj.conf}.
  \item Define the message type:
\begin{ccode}
struct battery_msg {
  uint8_t percent;
  bool charging;
};
\end{ccode}
  \item Define a channel named \code{battery_chan}, initially \code{\{.percent = 100, .charging = false\}},
    with a validator that rejects any message with \code{percent} above 100.
  \item Define a listener named \code{led}, observing the channel, whose callback prints \code{  led: red}
    (two spaces first) when the battery is below 20 \% and not charging, and \code{  led: green}
    otherwise.
  \item Write a thread \code{gauge}, priority \textbf{6}, stack 1024, no start delay, that publishes these
    readings, in this order:
\begin{ccode}
static const struct battery_msg readings[] = {
  {.percent = 80, .charging = false}, {.percent = 15, .charging = false},
  {.percent = 150, .charging = false}, {.percent = 15, .charging = true},
  {.percent = 60, .charging = true},
};
\end{ccode}

\noindent For each, it prints \code{gauge: <percent> \%}, followed by \code{, charging} if the
battery is charging, publishes the reading with a timeout of 50 ms, prints \code{gauge: rejected} if
the channel rejected it, and sleeps 100 ms.
\end{itemize}

The output must be:

\begin{console}
gauge: 80 %
  led: green
gauge: 15 %
  led: red
gauge: 150 %
gauge: rejected
gauge: 15 %, charging
  led: green
gauge: 60 %, charging
  led: green
\end{console}

\noindent \textbf{Hint:} A listener callback reads the message in place with
\code{zbus_chan_const_msg(chan)}. In a format string, \code{\%\%} prints a percent sign.

\task{Reflection}
\begin{itemize}
  \item In which thread does the \code{led} callback run? Check it: enable \code{CONFIG_THREAD_NAME=y}, and
    print \code{k_thread_name_get(k_current_get())} in the callback.
  \item What number does \code{zbus_chan_pub()} return for the 150 \% reading on \code{native_sim}? Why
    should the code compare it with \code{-ENOMSG} rather than with that number?
\end{itemize}

\exercise{A display that subscribes}{Code}
\label{c4:ex:3.2}

\task{Tasks}
\begin{itemize}
  \item Define a subscriber named \code{display}, with a notification queue of \textbf{4}, and add it to the
    channel’s observers after \code{led}.
  \item Write a thread \code{display_thread}, priority \textbf{7}, stack 1024, no start delay, which forever
    waits for a notification, reads the channel, and prints \code{  display: <percent> \%}, followed by
    \code{, charging} if the battery is charging.
\end{itemize}

The output must now be:

\begin{console}
gauge: 80 %
  led: green
  display: 80 %
gauge: 15 %
  led: red
  display: 15 %
gauge: 150 %
gauge: rejected
gauge: 15 %, charging
  led: green
  display: 15 %, charging
gauge: 60 %, charging
  led: green
  display: 60 %, charging
\end{console}

\task{Reflection}
\begin{itemize}
  \item Remove the gauge’s 100 ms sleep. What does the display print now, and why? Did it receive four
    notifications?
  \item Keep the sleep removed, and give \code{display_thread} priority 5, above the gauge. What does the
    display print now?
  \item Why does the \code{led} never miss a reading, whatever the priorities?
\end{itemize}

% ---------------------------------------------------------------------
\exerciseset{A garage door as a state machine}
\label{c4:set:4}
Create an application named \code{exercise4} to work in for this exercise set, with
\code{target_sources(app PRIVATE src/main.c src/door.c)} in its \code{CMakeLists.txt}.

\exercise{The state machine}{Code}
\label{c4:ex:4.1}
In this exercise you will write a hierarchical state machine with SMF, behind an API that tests can
drive directly.

\task{Tasks}
Copy this header into \code{src/door.h} unchanged. It is the interface the tests use:

\begin{ccode}
#ifndef DOOR_H_
#define DOOR_H_

#include <stdbool.h>

/** What can happen to the door. */
enum door_event {
	DOOR_EVENT_BUTTON,   /**< The button was pressed. */
	DOOR_EVENT_OPENED,   /**< The door reached its top end switch. */
	DOOR_EVENT_CLOSED,   /**< The door reached its bottom end switch. */
	DOOR_EVENT_OBSTACLE, /**< The light barrier saw something in the way. */
};

/** The states the door can be in. */
enum door_state {
	DOOR_STATE_CLOSED,
	DOOR_STATE_OPENING,
	DOOR_STATE_OPEN,
	DOOR_STATE_CLOSING,
	DOOR_STATE_STOPPED,
};

/** @brief Start the state machine, in the closed state, with the motor off
 *         and never started. */
void door_init(void);

/** @brief Run the state machine once, for one event. */
void door_handle(enum door_event event);

/** @brief Return the state the door is in. */
enum door_state door_get_state(void);

/** @brief Return whether the motor is running. */
bool door_motor_on(void);

/** @brief Return how many times the motor has been started since door_init(). */
unsigned int door_motor_starts(void);

#endif /* DOOR_H_ */
\end{ccode}

\noindent In \code{prj.conf}, enable SMF with parent states. In \code{src/door.c}, implement the
header with SMF, with these six states:

\begin{booktable}{@{}llp{9.5em}L@{}}
  \thead{State} & \thead{Parent} & \thead{Entry} & \thead{Handles} \\
  \midrule
  \code{MOVING} & none & Motor on, and count the start & \code{BUTTON}: go to \code{STOPPED}. Every other event: ignore it. Exit: motor off. \\
  \code{CLOSED} & none &  & \code{BUTTON}: go to \code{OPENING}. \\
  \code{OPENING} & \code{MOVING} &  & \code{OPENED}: go to \code{OPEN}. Anything else: pass it to the parent. \\
  \code{OPEN} & none &  & \code{BUTTON}: go to \code{CLOSING}. \\
  \code{CLOSING} & \code{MOVING} &  & \code{CLOSED}: go to \code{CLOSED}. \code{OBSTACLE}: go to \code{OPENING}. Anything else: pass it to the parent. \\
  \code{STOPPED} & none &  & \code{BUTTON}: go to \code{CLOSING}. \\
\end{booktable}

Events a state does not list are ignored. \code{door_get_state()} returns the state the door is in,
and \code{MOVING} is never one: the door is always in one of the five states of the header.

\textbf{Hint:} \code{MOVING}’s entry and exit actions are the only code that touches the motor. Keep
the state in a variable that each leaf state’s entry action sets.

\exercise{Drive it}{Code}
\label{c4:ex:4.2}

\task{Tasks}
In \code{src/main.c}, call \code{door_init()}, print the door’s status as \code{start}, and then
handle these events one at a time, printing the status after each, labelled with the event’s name:

\begin{ccode}
static const enum door_event script[] = {
	DOOR_EVENT_BUTTON,   DOOR_EVENT_OBSTACLE, DOOR_EVENT_OPENED,
	DOOR_EVENT_BUTTON,   DOOR_EVENT_OBSTACLE, DOOR_EVENT_OPENED,
	DOOR_EVENT_BUTTON,   DOOR_EVENT_BUTTON,   DOOR_EVENT_BUTTON,
	DOOR_EVENT_CLOSED,   DOOR_EVENT_OBSTACLE,
};
\end{ccode}

\noindent Print each status with \code{"\%-8s -> \%-7s motor \%-3s (started \%u times)\\n"}: the
label, the state’s name, \code{on} or \code{off}, and the number of motor starts. The output must
be:

\begin{console}
start    -> CLOSED  motor off (started 0 times)
BUTTON   -> OPENING motor on  (started 1 times)
OBSTACLE -> OPENING motor on  (started 1 times)
OPENED   -> OPEN    motor off (started 1 times)
BUTTON   -> CLOSING motor on  (started 2 times)
OBSTACLE -> OPENING motor on  (started 2 times)
OPENED   -> OPEN    motor off (started 2 times)
BUTTON   -> CLOSING motor on  (started 3 times)
BUTTON   -> STOPPED motor off (started 3 times)
BUTTON   -> CLOSING motor on  (started 4 times)
CLOSED   -> CLOSED  motor off (started 4 times)
OBSTACLE -> CLOSED  motor off (started 4 times)
\end{console}

\noindent The tests also drive \code{door.c} directly, without \code{main.c}, through sequences of
events the script does not contain.

\task{Reflection}
\begin{itemize}
  \item \code{OBSTACLE} while closing sends the door back to \code{OPENING}, and the motor start count stays
    at 2. Which rule of SMF makes that happen?
  \item Which state handles \code{OBSTACLE} while the door is opening, and what does it do with it?
  \item Change \code{OPENING}’s run action to return \code{SMF_EVENT_HANDLED} for every event. What happens
    when the button is pressed while the door is opening? Which tests fail?
  \item Remove \code{CONFIG_SMF_ANCESTOR_SUPPORT=y} from \code{prj.conf}. Does it still build? What does the
    program print now, and why?
  \item Why can \code{door.c} be tested without threads, timers, or \code{native_sim}’s clock?
\end{itemize}
